EXERCISE (cf. Exercise 1.4.E). The purpose of this exercise is to give you more practice with “left adjoints of forgetful functors,” the means by which we get abelian groups from abelian semigroups, and sheaves from presheaves. Suppose $A$ is a ring, and $S$ is a multiplicative subset. Then $S^{-1}A$-modules are a full subcategory (§1.1.15) of the category of $A$-modules (via the obvious inclusion $\operatorname{Mod}_{S^{-1}A}\hookrightarrow\operatorname{Mod}_A$). Then $\operatorname{Mod}_A\to\operatorname{Mod}_{S^{-1}A}$ can be interpreted as an adjoint to the forgetful functor $\operatorname{Mod}_{S^{-1}A}\to\operatorname{Mod}_A$. State and prove the correct statements. (Here is the larger story. Every $S^{-1}A$-module is an $A$-module, and this is an injective map, so we have a forgetful functor $F:\operatorname{Mod}_{S^{-1}A}\to\operatorname{Mod}_A$. In fact this is a fully faithful functor: it is injective on objects, and the morphisms between any two $S^{-1}A$-modules as $A$-modules are just the same when they are considered as $S^{-1}A$-modules. Then there is a functor $G:\operatorname{Mod}_A\to\operatorname{Mod}_{S^{-1}A}$, which might reasonably be called “localization with respect to $S$,” which is left-adjoint to the forgetful functor. Translation: If $M$ is an $A$-module, and $N$ is an $S^{-1}A$-module, then $\operatorname{Mor}(GM,N)$ (morphisms as $S^{-1}A$-modules, which are the same as morphisms as $A$-modules) are in natural bijection with $\operatorname{Mor}(M,FN)$ (morphisms as $A$-modules).)
